\scanpage{22}
\mspage{11}
a transitive system of isomorphisms between the $\pi_t$ $(t\in\unsure{\tau|\mathcal T})$, allowing them all to be identified with the same $T$. Thus $\langle\unsure{\tau|\mathcal T},\pi\rangle$ becomes a $T$-groupoid, that is, a groupoid (not necessarily connected) where the fundamental groups of the connected components are commutative and all identified with $T$.
Moreover, $\unsure{\mathcal E|\Sigma}$ acts on this $T$-groupoid (also acting on $T$) via its actions on the system $(\unsure{\tau|\mathcal T},\pi)$ (left $\pi$-space) by acting on $\unsure{\tau|\mathcal T}$ by the given action and on $\pi$ via inner automorphisms of $\unsure{\mathcal E|\Sigma}$, finally on $T$ via $\mathcal G$ via $\chi\ldots.$

[If $\unsure{\mathcal E|\Sigma}$ acts on $\unsure{\tau|\mathcal T}$, restricting this action to $\pi$ one finds a groupoid $\langle\unsure{\tau|\mathcal T},\pi\rangle$ on which $\unsure{\mathcal E|\Sigma}$ acts. Assuming this groupoid abelian, one thus has a family of abelian groups \inserted{$(T_j)_{j\in J}$} indexed by $J=\unsure{\tau|\mathcal T}/\pi$, on which $\unsure{\mathcal E|\Sigma}$ acts via $\unsure{\mathcal E|\Sigma}/\pi=\mathcal G$. This said, one assumes given a trivialization of this family $(T_j)$ by $T\simeq T_j$, compatible with action of
\mspage{21}
